\section{Methodology}\label{sec:methodology}

\subsection{Design: primitives first}\label{sec:method-primitives}

Throughout, ``derived'' means generated by the stated scenario mapping, not causally estimated. The central demand parameter is a transparent unit stress, $\varepsilon=1.0$, rather than an estimate. The grid includes $\varepsilon=0.4$ as an OBR-style lower anchor, 2.0 as the former April-outturn high case, and 3.0 as a severe sensitivity. The April movement is not used to identify the centre because it also reflects anticipation, front-running, prices, exchange rates, composition and other demand conditions. Likewise, $\pi=1.0$ is a constant-wage-share benchmark: the wage bill falls in proportion to the calibrated demand loss. The parameter grid is scenario sensitivity, not a statistical confidence interval.

A reader is entitled to ask, before any number is reported, which parts of this design are read off data and which are assumed. The answer, stated in full:

\begin{center}
\small
\begin{tabular}{p{0.28\textwidth}p{0.62\textwidth}}
\toprule
Object & Status \\
\midrule
Tariff schedule $\tau_j$ & \textbf{Data} (policy documents: the April 2025 schedule, the May 2025 EPD, the December 2025 pharmaceutical agreement). \\
US-export intensity $x_j$ & \textbf{Data} (HMRC RTS/OTS customs exports over ONS ABS division turnover; Appendix~\ref{app:support}). \\
Export fall & \textbf{Data} in the measured family (realised HMRC outturn, May 2025--February 2026); \textbf{assumed} in the tariff families through declared low, central, high and severe demand scenarios. \\
Pass-through $\pi$ & \textbf{Assumed} ($\pi = 1.0$), varied on a grid (Appendix~\ref{app:grid}). \\
Adjustment margin & \textbf{Assumed}, and deliberately varied. This is the paper's free dimension. \\
Household incidence & \textbf{Computed} from statutory tax and benefit rules in PolicyEngine UK; no behavioural or reduced-form step. \\
\bottomrule
\end{tabular}
\end{center}

The margin is assumed because the literature does not agree on it, and the disagreement is structural rather than empirical: \citet{ignatenko2025} derive the employment response from elastic labour supply and flexible wages, while \citet{rodriguezclare2025} derive it from downward nominal wage rigidity, so the same tariff primitive arrives as a wage shock in the first model and a job shock in the second. I do not endogenise the margin. I price the disagreement---holding the aggregate earnings loss fixed and reporting what each pole implies for households and the Exchequer. That is a weaker claim than a structural model makes, and a more auditable one.

The shock enters through primitives, not through an assumed employment effect. The pipeline is:
\begin{equation}
s_j \;=\; \tau_j \, x_j \, \varepsilon \, \pi,
\label{eq:shock}
\end{equation}
where $s_j$ is an imposed employee-wage-bill stress for SIC 2007 division $j$, $\tau_j$ is its tariff rate, $x_j$ is US-export intensity, $\varepsilon$ is an export-demand calibration and $\pi$ is a \emph{wage-bill-incidence assumption}. Calling $\pi$ pass-through would conflate distinct mechanisms. A structural transmission would separately map the tariff into export prices and quantities, quantities into domestic production and value added, and production into factor demand, productivity, profits, wages, hours and employment. Equation~\eqref{eq:shock} compresses those links into one accounting bridge. It neither estimates a production function nor identifies the causal response of productivity, output or labour demand.

$x_j$ is US goods exports as a share of division turnover, built from HMRC trade data and ONS Annual Business Survey turnover denominators; Appendix~\ref{app:support} records the frozen provenance and the turnover denominator used. The central $\varepsilon=1.0$ is a transparent unit stress, not an elasticity estimate; it amounts to a unit tariff-to-export-demand semi-elasticity, under which each percentage point of tariff removes one per cent of US-bound demand, and $\varepsilon=0.4$ is the lower anchor, in the spirit of the trade-elasticity assumptions underlying the OBR's March 2025 tariff scenarios \citep{obr2025}. The central $\pi=1$ scales each division's wage bill in proportion to its calibrated demand loss---a constant-wage-share benchmark, not full incidence, which would put the entire revenue loss on the wage bill and require a $\pi$ several times larger; the benchmark is an assumption, not a literature estimate. The sensitivity grid varies $\varepsilon\pi$, but this is parameter sensitivity, not an uncertainty interval around an identified coefficient. The implied gross earnings loss is therefore a conditional scenario input. OECD METRO and the UK Trade Modelling Review instead place tariff shocks inside multi-sector production systems in which prices, output, value added and factor demand adjust jointly \citep{luu2020, ditreview2021}; the World Bank HIT framework similarly labels its direct household calculations first-order and excludes productivity and substitution responses \citep{artuc2021}. This paper begins only after those missing production-side mechanisms would have generated a labour-income distribution.

\subsection{Public product--destination benchmark}\label{sec:method-trade-panel}

To test whether post-policy US exports move differently from other UK export
destinations, I build a monthly HMRC OTS panel from January 2018 to May 2026
for the United States, Canada, Japan and Australia \citep{hmrcots2026}.
Detailed HS8 rows are aggregated to HS4 and balanced across products,
destinations and months, with absent flows set to zero. HMRC's separate
two-digit aggregate rows are excluded. For product $p$ and month $t$, the
outcome is
\[
g_{pt}=\log(1+V_{p,\mathrm{US},t})
-\frac{1}{3}\sum_{d\in\{\mathrm{CA,JP,AU}\}}\log(1+V_{pdt}).
\]
I regress this destination gap on a post-May-2025 indicator, HS4 product
fixed effects, month-of-year effects and a linear trend. April 2025 is
omitted as an announcement and anticipation month. The primary WLS weights
are products' mean 2022--24 US export values, capped at the 99th percentile;
standard errors are clustered by HS4. Uncapped value weights, equal product
weights, individual destinations, later sample starts and pseudo-policy dates
in May 2019--23 are reported as specification checks.

This design controls for product shocks common across destinations but not
US-specific demand, exchange rates, shipping, inventories, composition or
other coincident policies. The high-tariff comparison groups HS chapters 72,
73 and 87 as steel and automotive products; that broad mapping is a
falsification check, not a complete product-level statutory tariff schedule.
The panel is therefore not used to estimate $\varepsilon$ or $\pi$ in
Equation~\eqref{eq:shock}. It tests whether the imposed bridge is consistent
with destination-differential outturn patterns.

\subsection{The tariff schedule and the EPD counterfactual}\label{sec:method-schedule}

Two tariff scenarios bracket the policy. \emph{Full tariffs}: the April 2025 schedule---10 per cent baseline on UK goods, 25 per cent on motor vehicles (SIC 29) and steel (SIC 24). \emph{EPD}: the deal-mitigated schedule---autos at the 10 per cent in-quota rate on up to 100,000 vehicles; steel with conditional, partial relief; pharmaceuticals (SIC 21) exempt under the December 2025 agreement. The automotive scenario applies the 10 per cent rate to the modelled 2024 exposure because the quota is approximately equal to that year's US-bound volume; it does not model quota allocation, above-quota growth or an endogenous volume response around the ceiling. The difference between the two scenarios, run through identical downstream machinery, prices the EPD for households and the Exchequer. Two caveats travel with the schedule. Steel: about 80 per cent of UK steel exports go to the EU, so the US tariff is a minority driver of the sector's exposure and the steel cell should be read alongside EU/UK safeguards. Pharmaceuticals: the exemption was purchased through VPAG rebate cuts and NHS pricing concessions, so pharma's incidence is fiscal and price-side, not an employment shock; I model its employment shock as zero under the EPD and discuss the fiscal channel separately (Section~\ref{sec:policy}).

\subsection{Adjustment-margin families}\label{sec:method-margins}

The derived $s_j$ is realised on Family Resources Survey 2024--25 employees
in exposed divisions. The submission design treats the mixed adjustment
surface as the conceptual centre: statutory cushioning is evaluated for a
common expected sector wage-bill loss while the share delivered through
temporary non-employment, survivor earnings reductions, and re-employment is
varied. The pure displacement and proportional wage-cut cases are transparent
bounds on that surface. Inactivity, literal sectoral reallocation and the EPD
schedule are complementary sensitivities, not alternative causal estimates.
Wage cuts remove $\sum_j s_j B_j$ deterministically; displacement removes
that amount in expectation under the risk model and is numerically integrated
using the balanced design below.

\paragraph{Displacement (ADH short run).} Within each exposed division every
employee has baseline risk $s_j$, irrespective of the record's survey weight.
The primary estimator generates repeated random-start systematic candidates,
ordered across age and within-industry earnings ranks, and retains the
candidate that best matches the expected industry wage-bill and weighted-
headcount losses while secondarily balancing broad age and earnings margins.
This is balanced numerical integration, not a survey estimator: conditioning
on aggregate balance intentionally changes exact record-level marginal
inclusion probabilities, and because the retained candidate minimises a
score dominated by the total wage-bill target, per-division losses are
matched softly rather than exactly and best-of-candidates selection carries
no formal unbiasedness guarantee. The claim that the balanced design removes
the declared amount in expectation is therefore established empirically, by
its agreement with the unbiased comparator, not by construction. Selected
workers move to unemployment with earnings, hours, pension contributions and
statutory pay zeroed. An independent Bernoulli estimator, which preserves
the declared first-order probabilities exactly but has much larger
allocation dispersion, is reported as the unbiased assignment-rule
comparator; at the unit scale the two agree on cushioning to half a
percentage point (Section~\ref{sec:results-cushioning}).
\paragraph{Wage cuts with reallocation (Dauth/Traiberman long run).} No job loss; every employee in division $j$ takes a proportional cut of $s_j$, so the weighted aggregate earnings removed equals $\sum_j s_j B_j$ by construction---exact earnings-equivalence with the displacement family's expected removal. Employee pension and salary-sacrifice contributions are scaled proportionally with the cut, mirroring their zeroing under displacement; this lowers taxable pay (and hence tax relief) but also reduces the household's pension outflow, so under the HBAI disposable-income concept part of the earnings loss is absorbed by reduced retirement saving rather than by current income.
\paragraph{Inactivity (Beatty--Fothergill UK history).} The displacement draw, but displaced workers aged 50 and over exit to economic inactivity, the benefit-financed route of UK deindustrialisation, rather than unemployment. The route is benefit-financed in the model, not just relabelled: every inactive worker is flagged as having limited capability for work-related activity, so their household's Universal Credit includes the LCWRA health element (\pounds 217.26 a month in 2026---the reduced rate applying to new claims from April 2026 under the Universal Credit Act 2025; existing claimants retain the higher pre-reform rate, and all modelled inactivity claims are new claims). In policyengine-uk 2.89.2 the survey-imputed \texttt{uc\_LCWRA\_element} is stored, so the post-shock benefit-unit value is set to the maximum of the stored amount and one annual statutory element; an existing recipient is never paid a second element, unaffected units are unchanged, and the run hard-errors if the element fails to reach a newly flagged person's benefit unit. This mirrors the \citet{beatty2007} route by which industrial job loss became incapacity-benefit inactivity. The strong assumption should be stated plainly: \emph{all} displaced workers aged 50 and over are assumed to pass the Work Capability Assessment, so the family is an upper bound on the health-element route; no partial-qualification sensitivity (for example, 50 per cent of the older displaced qualifying) has been simulated, which is a limitation of the present scenario set. Work-search conditionality carries no separate fiscal machinery in the model, so the LCWRA element is the entire modelled wedge between inactivity and unemployment.

\paragraph{Reallocation into services (ADH/\citealp{autor2014}).} The same displacement draw on the same seed---so the two families are paired worker-for-worker---but the drawn workers are re-employed rather than left out of work. Their \texttt{sic\_industry\_division} is reassigned to one of four services destinations (SIC 47 retail, 86 human health, 49 land transport, 56 food and beverage service) in proportion to those divisions' FRS employee-weight shares (29.5, 45.8, 9.2 and 15.5 per cent), and their annual earnings are cut by a penalty calibrated on the FRS itself: weighted mean annual employee earnings of \pounds 48{,}272 over 2{,}047 hours in the exposed goods divisions against \pounds 34{,}594 over 1{,}701 hours in the four destinations, a \textbf{28.3 per cent annual-earnings penalty} that embeds the hours fall. Two variants are reported: an hours- and age-controlled \textbf{low-penalty} case of 14.0 per cent (the pure hourly-wage gap), and a \textbf{three-month lag} in which workers are out of work for a quarter of the year before the services job starts, scaling annual earnings and hours by a further $1-3/12$. ASHE 2024 full-time medians bracket the hours-controlled figure at 10--20 per cent; the FRS all-employee penalty is larger because it also captures the shift into part-time work. Because this family removes only the penalty on the movers rather than the full earnings-equivalent aggregate, its pound figures are comparable with the displacement family (identical worker sets, paired draws) but \emph{not} with the wage-cut family; its cushioning \emph{rate} is comparable with both. Appendix~\ref{app:support} reports the results and mechanics.

\paragraph{Common-loss transition path.} The redesigned transition margin
combines a transition-equivalent annual earnings/hours penalty, service-sector
re-employment and survivor earnings reductions while preserving the common-loss
normalization. Let $s$ be
the sector shock, $\lambda$ the share of the declared loss assigned to
transition workers, and $d=1-(1-r)(1-\ell/12)$ the annual earnings loss of a
transition worker, where $r$ is the re-employment penalty and $\ell$ the lag in
months. Transition probability is $p=\lambda s/d$; survivors take a cut
$c=(1-\lambda)s/(1-p)$. Hence $pd+(1-p)c=s$ in expectation. The illustrative
interior calibration uses $\lambda=0.70$, $r=0.283$ and $\ell=3$; the
transition share is motivated by the displacement-cost decompositions of
\citet{rud2022} and \citet{dixcarneiro2014}, which attribute the majority of
displacement costs to non-employment spells and re-employment penalties.
But those studies observe only displaced workers and cannot bound the
survivor share of an aggregate wage-bill loss, so $\lambda=0.70$ is an
assumption, not a literature-identified quantity. It is
one point on the response surface, which remains the paper's central
object. In the annual
simulation, the lag is represented by an earnings/hours factor: workers remain
employed for PolicyEngine purposes, so no within-year unemployment or JSA/UC
spell is generated. This is an accounting normalization for comparing paths,
not an estimate of the probability of re-employment or the UK tariff response.
The 14 per cent hours/age-controlled
penalty and zero-, three- and six-month lags are sensitivity dimensions.

\paragraph{Mixed adjustment and factorial scenario testing.} Let $\lambda\in[0,1]$ denote the share of the sector shock delivered through temporary non-employment rather than an incumbent earnings reduction. A worker with sector shock $s$ is displaced with probability $p=\lambda s$. Conditional on surviving, the proportional wage cut is
\[
c=\frac{s-p}{1-p},
\]
so expected earnings loss is exactly $p+(1-p)c=s$ for every worker. The endpoints reproduce pure wage cuts ($\lambda=0$) and pure temporary displacement ($\lambda=1$); because the mixed family draws displacement independently, the $\lambda=1$ endpoint reproduces the Bernoulli comparator rather than the balanced primary estimator. A separate reallocation family moves displaced workers into observed service destinations, with either an immediate earnings penalty or a three-month non-employment lag. It is paired worker-for-worker with displacement but does not preserve the common aggregate loss unless its penalty is explicitly re-scaled, so its pound totals are reported as a transition sensitivity rather than combined with the mixed grid. The exploratory grid crosses $\lambda\in\{0,.25,.5,.75,1\}$ with export-demand calibrations $\varepsilon\in\{0.4,1,2,3\}$ on five common assignments per cell. This follows the companion UK AI study's factorial principle: vary shock magnitude separately from adjustment composition. It is scenario testing, not a probability distribution over future outcomes.

\paragraph{The wage-cut margin as an upper bound, and a literature-disciplined mixed sensitivity.} The pure wage-cut family should be read against the rent-sharing evidence, because it embeds a survivor-wage response far above anything estimated. Delivering the entire sector wage-bill loss through proportional cuts to incumbent workers' wages imposes a unit wage response to the sector exposure index. \citet{garinsilverio2024} estimate an incumbent-wage elasticity of roughly 0.13--0.14 to firm-specific output shocks in Portuguese matched data; the rent-sharing survey of \citet{cardcardosoheining2018} places typical elasticities at 0.05--0.15; and \citet{hummels2014} find small within-job wage responses in Danish data alongside losses through displacement. The wage-cut family ($\lambda=0$) is therefore an \emph{upper-bound stress scenario for survivor-wage incidence}, retained because it prices the flexible-wage pole of the structural disagreement.

The scenario surface also contains a 15/85 point ($\lambda=0.85$): 15 per cent of the imposed wage-bill loss delivered through survivor wage cuts, disciplined by the upper end of the cited wage-response estimates. A wage elasticity to firm output or rents is not algebraically the share of an aggregate wage-bill loss borne by survivor wages---identifying that share requires linked employer--worker data---so this is one surface point among others, subordinate to the transition path above as the paper's single illustrative interior calibration.

Two scope limits on the family set should be stated here rather than in a footnote. The wage-cut family models the \emph{earnings consequence} of reallocation without moving anyone between sectors: workers keep their observed industry and hours, and the proportional cut stands in for the wage penalty that a completed reallocation would carry; the reallocation family does move them, but who reallocates is assumed (it is the displacement draw) rather than estimated. And no household-side behavioural margin is active---added-worker responses, earlier or later retirement, and shifts into self-employment \citep{irastorza2025} are all suppressed, as are any labour-supply responses to the tax and benefit changes themselves. All four families are therefore first-round, rules-based incidence.

A third caveat concerns selection within divisions. Export-exposed workers are
unlikely to be representative of their division: exporters pay a wage premium
and pass export-demand shocks into incumbent wages
\citep{hummels2014, garinsilverio2024}. Cushioning is nonlinear in earnings,
so uniform risk may misstate its tax/benefit composition. The longitudinal-LFS
cell, earnings-band and QRF sensitivities vary observed risk shape while
recalibrating each division to its wage-bill target. They quantify selection-
model sensitivity but cannot identify exporter status or tariff-specific
worker risk.

The FRS support for the transition is thin. Of 34,966 person records, 12,800
have employee earnings and 1,018 are in mapped exposed divisions, but the
reference probabilities imply only 9.7 selected records. The largest exposed
record has a person weight above 19 thousand. Balanced integration constrains
the realised aggregate and observed composition; it does not create survey
information or establish population representativeness. I therefore report
the affected record count, Kish-style effective number of gross-loss
contributions, largest-record loss share, and balanced-versus-Bernoulli mean
and SD comparisons. Subgroup rankings remain exploratory.

\paragraph{Duration-equivalent annual stresses.} The annual PolicyEngine model
cannot represent a worker employed for part of a year and unemployed for the
remainder with internally consistent monthly UC assessment, work status and
earnings. I therefore do not simulate three- or six-month unemployment
spells. Instead I scale $s_j$ by $d/12$, for
$d\in\{3,6,12\}$, and assign a correspondingly smaller number of coherent
full-year displaced states. These rows preserve the expected worker-month
earnings loss associated with each duration while changing the extensive-
margin stock. They are labelled \emph{duration-equivalent annual stresses},
not spell estimates. Wage-cut rows use the same scale, preserving the common-
loss comparison at every duration and at both the OBR-style 0.4 and unit
anchors. Appendix~\ref{app:monthlyuc} bounds the bias from this annualisation
with a parameter-based comparison of the duration-equivalent stress against a
correct monthly UC assessment for a representative displaced worker: UC
entitlement per worker-month is identical by construction, the primary
12-month contrast is unaffected, and the residual differences at 3- and
6-month scales run through partial-year income tax and payment timing.

The margin is the paper's central axis because the tax--benefit system treats the two poles asymmetrically: an in-work earnings reduction is offset automatically at marginal tax and taper rates, while a job loss moves the worker to the gated, means-tested extensive margin of the benefit system. How much each pole is cushioned, and---the question that turns out to discriminate between them---how much that cushioning depends on anyone doing anything, are empirical questions the microdata answer (Section~\ref{sec:results-cushioning}).

\subsection{Longitudinal LFS imputation benchmark}\label{sec:method-lfs}

The uniform within-division assignment in the central displacement family is
transparent but suppresses observable worker heterogeneity. I therefore use
UK Data Service Study 9490, the Labour Force Survey Five-Quarter Longitudinal
Dataset covering July 2024--September 2025, as a donor benchmark. The donor
population comprises respondents employed at wave 1. Job exit is observed at
wave 5; log wage change is defined only for respondents employed with
positive main-job pay at both endpoints. Negative UKDS sentinel values are
missing, not zero. Public 2024 Nomis BRES employee totals rescale the
manufacturing SIC composition before direct transition targets are computed.
\citep{onslfs9490}

The primary imputation maps SIC-division--sex--age cells to employed FRS
receivers, shrinking sparse cells towards sector and manufacturing means.
After matching, predicted exit risks are rescaled so their FRS-weighted mean
equals the directly measured, BRES-composition-adjusted LFS manufacturing
exit rate. Mean predicted log wage changes are centred on the corresponding
LFS target. An income-tercile estimator supplies a transparent sensitivity
specification.

A regularised quantile/random-forest benchmark follows the calibrated
imputation architecture in PolicyEngine's \texttt{young-worker-nics}
project. It estimates model shape on all wave-1 employed adult donors using
age, sex and annualised employment income. The binary outcome uses the
random-forest event probability; the continuous outcome uses the QRF
conditional mean. Forests use 500 trees and a minimum leaf size of 20; wage
outcomes are winsorised at the donor 1st and 99th percentiles. Both receiver
predictions are then calibrated to exactly the same manufacturing LFS
targets as the cell estimator. This is a predictive baseline-transition
benchmark, not an estimate of tariff-induced exits or wage changes. It is
therefore reported as a robustness diagnostic and does not replace the
declared central assignment rule.

In a downstream allocation sensitivity, each receiver score $r_i$ is
converted to a within-division displacement probability
$p_i=\min(1,k_jr_i)$, where $k_j$ is solved numerically so
\[
\frac{\sum_{i\in j}p_i E_iw_i}{\sum_{i\in j}E_iw_i}=s_j.
\]
Thus cells, earnings bands and QRF change who is exposed while preserving
each division's expected employee-wage-bill loss. Missing scores use the
division mean. Common assignments then isolate sensitivity to incidence
shape; they do not give the baseline transition model a causal tariff
interpretation.

\subsection{Universal Credit take-up after the shock}\label{sec:method-takeup}

The take-up experiment is applied symmetrically across margins. After each shock, changed benefit units whose all-claim UC award moves from zero at baseline to positive post-shock receive a claiming draw at the scenario take-up rate. This includes wage-cut households that cross into entitlement. Existing claimants, existing eligible non-claimants, unchanged units and post-shock-ineligible units retain their baseline flag, avoiding an unintended population-wide take-up reform. The grid therefore measures sensitivity to a common new-entitlement convention; it does not estimate actual claiming behaviour.

One modelling choice deserves its own subsection, because an earlier version of this paper got it wrong; the correction turns out to leave the unit-stress headline unchanged (Section~\ref{sec:results-cushioning}), but that is a measured result, not a licence for the incorrect convention. PolicyEngine UK carries benefit take-up as a benefit-unit-level \emph{stored input}, \texttt{would\_claim\_uc}, drawn once at dataset build time. Two properties of that draw matter. First, its target is a flat population-wide parameter---0.55 as a share of \emph{all} benefit units, with reported UC recipients anchored to \textsc{true} and the remainder filled to hit the target. It is not a take-up rate among the entitled, and it is therefore not comparable to the roughly 80 per cent take-up-among-entitled that the Resolution Foundation and the IFS assume; in the dataset used here the realised weighted flag rate is 0.547 over all benefit units and 0.541 among UC-eligible ones. Second, the flag is drawn on \emph{pre-shock} circumstances. Nothing in a naive shock pipeline re-draws it, so a worker displaced by the shock carries into unemployment the claiming decision attributed to them while they were employed and, in the great majority of cases, not entitled to anything---measured take-up among the displaced post-shock was 46.9 per cent, below even the population flag rate. The baseline flag is thus not informative about the post-shock claiming behaviour of a newly entitled family: it encodes a decision taken in a state of the world the shock has destroyed.

The pipeline therefore re-draws \texttt{would\_claim\_uc} after each shock only for benefit units whose circumstances changed and whose all-claim UC award moves from zero at baseline to positive post-shock, at a scenario parameter \texttt{uc\_takeup} with default 0.80. The re-draw uses an independent random stream, leaving worker assignment unchanged and preserving paired comparisons. This rule applies to displacement, inactivity, reallocation and wage cuts alike; existing eligible, post-shock-ineligible and unchanged units retain their baseline flag. Because the entire displacement-side cushion runs through UC, take-up is treated as a headline sensitivity rather than an exploratory one: Section~\ref{sec:results-cushioning} reports the primary displacement margin at new-entitlement take-up rates of 0.55, 0.80 and 1.00 over the full comparator draw count. The 0.80 default is a strong assumption for a population with no UC claiming history and unobserved capital (redundancy payments would gate many claims under the frozen capital thresholds); the grid brackets it rather than estimating it.

\subsection{Microsimulation and outcomes}\label{sec:method-microsim}

The Monte Carlo draws describe sensitivity to the artificial assignment of displacement across weighted FRS records. Their standard deviations are assignment dispersion, not sampling standard errors, confidence intervals, or uncertainty about the tariff response, survey design or model parameters. Numerical Monte Carlo error of a reported mean is the assignment SD divided by the square root of the number of draws. Parameter scenarios vary the demand calibration, wage-bill incidence, take-up and adjustment composition, but are not assigned probabilities and do not form confidence intervals. Survey sampling variance is approximated for the primary contrast by a household bootstrap of the FRS sample (Section~\ref{sec:results-cushioning}); the research file carries no PSU or stratum identifiers, so the bootstrap treats households as independent and cannot reflect geographic clustering. Uncertainty in the trade-to-earnings bridge is not estimated. Reported comparisons other than the bootstrapped primary contrast are descriptive scenario contrasts. Distributional cash-income changes use annual household disposable-income change divided by household size before person-weighted aggregation; aggregate disposable-income totals remain household-weighted totals.

The realisation step---an imposed status and earnings transition on survey households, then run through the full benefit rules---follows \citet{brewertasseva2021}, who apply it to the Covid-19 labour-market shock in UKMOD, and the EUROMOD nowcasting tradition of imposed labour-market transitions \citep{navicke2013}. PolicyEngine UK computes baseline and shocked household disposable income (HBAI cash concept throughout), relative before-housing-costs poverty measured against a poverty line frozen at 60 per cent of the baseline equivalised median (so the reported change reflects households crossing a fixed threshold rather than mechanical movement of the line itself), absolute before-housing-costs poverty against a fixed baseline line, after-housing-costs poverty, the Gini coefficient, distributional breakdowns and the Exchequer effect. The \emph{Exchequer effect} is defined as the change in the aggregate simulated government balance---all modelled tax and contribution revenue (including employer National Insurance contributions) minus all modelled benefit spending---between the baseline and shocked simulations. It is therefore broader than the employee-earnings loss $\Delta E$ that enters the cushioning rate; Appendix~\ref{app:support} reconciles the two. On timing: trade exposure is built from 2024 HMRC customs data, the household input is FRS 2024--25, and the simulation year is 2026, with PolicyEngine UK's standard uprating applied to survey monetary variables between the survey year and the simulation year. Primary submission scenarios use \SubmissionDraws\ balanced assignments; legacy direct and measured scenarios use \ProductionDraws\ assignments, and explicitly exploratory diagnostics use five. Following \citet{ignatenko2025}, revenue recycling can flip aggregate employment responses; here fiscal policy is held fixed. Only direct taxes and benefits respond. Indirect taxes are not shocked, so reduced consumption-tax receipts are absent from the government balance.
